% chapters/03_methodik.tex
% Drittes Kapitel: Methodik
% Beschreibt das Forschungsdesign und die eingesetzten Methoden.

\chapter{Methodik}
\label{chap:methodik}

Dieses Kapitel beschreibt das methodische Vorgehen der Arbeit. Es wird
erlaeutert, wie die in Kapitel~\ref{chap:einleitung} formulierten
Forschungsfragen beantwortet werden.

\section{Forschungsdesign}
\label{sec:forschungsdesign}

Das Forschungsdesign legt den uebergeordneten Rahmen der Untersuchung fest.
Es wird unterschieden zwischen quantitativen, qualitativen und
Mixed-Methods-Ansaetzen \autocite{beispiel_buch_2020}.

\begin{itemize}
  \item \textbf{Quantitativ:} Messung und statistische Auswertung numerischer Daten.
  \item \textbf{Qualitativ:} Interpretation von Text, Interviews oder Beobachtungen.
  \item \textbf{Mixed Methods:} Kombination beider Ansaetze.
\end{itemize}

Diese Arbeit folgt einem [quantitativen / qualitativen / Mixed-Methods]-Ansatz,
da \ldots

\section{Datenerhebung}
\label{sec:datenerhebung}

Tabelle~\ref{tab:methoden} gibt einen Ueberblick ueber die eingesetzten
Erhebungsmethoden und deren Eignung fuer die jeweilige Forschungsfrage.

\begin{table}[h]
\centering
\caption{Uebersicht der Erhebungsmethoden}
\label{tab:methoden}
\begin{tabular}{lll}
\toprule
\textbf{Methode}    & \textbf{Forschungsfrage} & \textbf{Begründung} \\
\midrule
Methode 1           & FF1                      & Begruendung A       \\
Methode 2           & FF2                      & Begruendung B       \\
Methode 3           & FF1, FF3                 & Begruendung C       \\
\bottomrule
\end{tabular}
\end{table}

\subsection{Stichprobe und Auswahl}
\label{subsec:stichprobe}

Die Stichprobe umfasst \ldots Einheiten, die nach dem Kriterium \ldots
ausgewaehlt wurden. Die Auswahl erfolgte [zufaellig / bewusst / geschichtet],
weil \ldots

\subsection{Erhebungsinstrument}
\label{subsec:instrument}

Als Erhebungsinstrument wurde [Fragebogen / Interviewleitfaden / \ldots]
eingesetzt. Der Aufbau orientiert sich an \autocite{beispiel_artikel_2019}.

\section{Datenauswertung}
\label{sec:auswertung}

Die erhobenen Daten wurden mit [Software / Methode] ausgewertet.
Fuer quantitative Daten kamen deskriptive Statistiken sowie [Verfahren]
zum Einsatz. Qualitative Daten wurden nach der Methode der
[Inhaltsanalyse / Grounded Theory / \ldots] kodiert und kategorisiert.
